\documentclass[11pt]{article}
\usepackage[T1]{fontenc}
\usepackage[utf8]{inputenc}
\usepackage{amsmath,amssymb,amsthm,mathtools}
\usepackage{geometry,graphicx,booktabs,longtable,tabularx,array}
\usepackage{enumitem,needspace,xcolor,tikz,tcolorbox}
\usetikzlibrary{arrows.meta,positioning,calc,fit,backgrounds,decorations.pathreplacing,matrix,patterns}
\tcbuselibrary{breakable}
\usepackage{caption,hyperref,bookmark,cleveref,microtype}
\usepackage{komorebi-publication}
\renewcommand{\KomorebiTitle}{Komorebi Space}
\renewcommand{\KomorebiSubtitle}{Coherent realizations, obstructions, and descent}
\renewcommand{\KomorebiSeries}{Foundations of the Realization Calculus}
\renewcommand{\KomorebiKind}{Monograph}
\renewcommand{\KomorebiEdition}{}
\renewcommand{\KomorebiShortTitle}{Komorebi Space}
\renewcommand{\KomorebiProfile}{monograph}
\KomorebiSetup
\setcounter{tocdepth}{1}
\setcounter{secnumdepth}{2}
\setlist{nosep,leftmargin=1.7em,topsep=4pt}
\hypersetup{pdfkeywords={Komorebi Space,realization calculus,descent,obstruction,Lean,holonomy,closed range}}
\newtheorem{theorem}{Theorem}[section]
\newtheorem{proposition}[theorem]{Proposition}
\newtheorem{lemma}[theorem]{Lemma}
\newtheorem{corollary}[theorem]{Corollary}
\theoremstyle{definition}
\newtheorem{definition}[theorem]{Definition}
\newtheorem{example}[theorem]{Example}
\theoremstyle{remark}
\newtheorem{remark}[theorem]{Remark}
\numberwithin{equation}{section}
\newcommand{\R}{\mathbb R}
\newcommand{\Z}{\mathbb Z}
\newcommand{\Q}{\mathbb Q}
\newcommand{\N}{\mathbb N}
\newcommand{\K}{\mathbb K}
\newcommand{\V}{\mathcal V}
\newcommand{\W}{\mathcal W}
\newcommand{\id}{\operatorname{id}}
\newcommand{\im}{\operatorname{im}}
\newcommand{\ran}{\operatorname{ran}}
\newcommand{\Hom}{\operatorname{Hom}}
\newcommand{\Ob}{\operatorname{Ob}}
\newcommand{\Ext}{\operatorname{Ext}}
\newcommand{\Match}{\operatorname{Match}}
\newcommand{\Exact}{\operatorname{Exact}}
\newcommand{\dist}{\operatorname{dist}}
\newcommand{\Aut}{\operatorname{Aut}}
\newcommand{\Stab}{\operatorname{Stab}}
\newcommand{\Set}{\mathbf{Set}}
\newcommand{\Cat}{\mathbf{Cat}}
\newcommand{\norm}[1]{\left\lVert#1\right\rVert}
\newcommand{\abs}[1]{\left\lvert#1\right\rvert}
% Preserve source anchors without exposing editorial identifiers in print.
\newcommand{\source}[1]{\phantomsection\label{src:#1}}
\newcommand{\code}[1]{\texttt{\detokenize{#1}}}
\newcommand{\chap}[1]{\par\Needspace{10\baselineskip}\section{#1}}
\newcommand{\partnote}[2]{\clearpage\part{#1}\noindent\begin{minipage}{.92\linewidth}\large #2\end{minipage}\par\vspace{12pt}}
\newtcolorbox{insight}{breakable,colback=KomorebiSoft!55,colframe=KomorebiPine!50,boxrule=.45pt,arc=1pt,left=9pt,right=9pt,top=7pt,bottom=7pt}
\newtcolorbox{scopebox}{breakable,colback=KomorebiWarm,colframe=KomorebiOchre!55,boxrule=.4pt,arc=1pt,left=9pt,right=9pt,top=7pt,bottom=7pt}
\usepackage{float}
\newcommand{\fixedplate}[3]{\begin{figure}[H]\centering\resizebox{.96\linewidth}{!}{\input{figures/#1}}\caption{#2}\label{#3}\end{figure}}
\newcommand{\plate}[3]{\begin{figure}[!htbp]\centering\resizebox{.96\linewidth}{!}{\input{figures/#1}}\caption{#2}\label{#3}\end{figure}}
\tikzset{box/.style={draw=KomorebiPine!70,rounded corners=2pt,fill=KomorebiSoft!50,align=center,inner sep=8pt,font=\small},arr/.style={-{Latex[length=2.2mm]},draw=KomorebiPine,thick},thinarr/.style={-{Latex[length=1.8mm]},draw=KomorebiMuted},lab/.style={font=\small,align=center},dot/.style={circle,fill=KomorebiPine,inner sep=2.3pt}}
\renewcommand{\arraystretch}{1.18}
\raggedbottom
\widowpenalty=7000\clubpenalty=7000
\setlength{\emergencystretch}{3em}

% Compact only the contents entries, without reducing body type size.
\usepackage{etoolbox}
\makeatletter
\patchcmd{\l@section}{1.0em}{0.6em}{}{}
\makeatother

% Counts for the accepted formal source accompanying this publication.
% Fresh integrated verification totals for publication v1.15.0.
\newcommand{\EditionProductionModules}{162}
\newcommand{\EditionRoots}{4}
\newcommand{\EditionTargets}{166}
\newcommand{\EditionTheorems}{1,859}
\newcommand{\EditionDeclarations}{5,339}
\newcommand{\EditionControls}{56}
\newcommand{\EditionPositiveControls}{18}
\newcommand{\EditionNegativeControls}{35}
\newcommand{\EditionPolicyControls}{3}
\newcommand{\EditionVerifiedGroups}{16}
\newcommand{\EditionWrittenGroups}{0}
\newcommand{\EditionPrincipalGroups}{16}
\newcommand{\EditionSelectedMapEntries}{114}
\newcommand{\EditionEntrypoints}{\EditionRoots}
\newcommand{\EditionExplicitTheorems}{\EditionTheorems}
\newcommand{\EditionExportedDeclarations}{\EditionDeclarations}


% Keep readable half-page figures with their surrounding exposition.
\renewcommand{\topfraction}{.88}
\renewcommand{\bottomfraction}{.80}
\renewcommand{\textfraction}{.10}
\renewcommand{\floatpagefraction}{.78}
\setcounter{topnumber}{2}
\setcounter{bottomnumber}{2}
\setcounter{totalnumber}{3}

% Evidence scope is presented in a reader-facing appendix.
% Formal coverage is documented in chapters/11_written_extensions.tex.
% Editorial source labels remain in the manuscript for cross-references only.

\newcommand{\written}[1]{}

\renewcommand{\KomorebiKind}{Visual guide}
\renewcommand{\KomorebiSubtitle}{Sixteen visual investigations in coherent realization}
\hypersetup{pdftitle={Komorebi Space: Sixteen Visual Investigations},pdfsubject={Worked visual experiments in realization, correction, and limits}}
\setcounter{tocdepth}{1}
\newcommand{\scene}[4]{%
  \clearpage\phantomsection\addcontentsline{toc}{section}{#1\quad #2}%
  {\sffamily\small\bfseries\color{KomorebiMuted}VISUAL INVESTIGATION #1\par}%
  \vspace{4pt}{\sffamily\Large\bfseries\color{KomorebiPine}#2\par}%
  \vspace{8pt}#3\par\vspace{4pt}%
  \begin{center}\includegraphics[width=\linewidth,height=5.05in,keepaspectratio]{#4}\end{center}\vspace{2pt}}
\newcommand{\scenenotes}[3]{%
  \noindent\textbf{Read the picture.} #1\par\medskip
  \noindent\textbf{Change one thing.} #2\par
  \vfill{\small\color{KomorebiMuted}Definitions, proofs, and general hypotheses: #3.\par}}
\begin{document}
\KomorebiCover
\clearpage
\thispagestyle{empty}
\section*{Publication and reuse}
\noindent\emph{Komorebi Space: Coherent Realizations, Obstructions, and Descent}\\
Komorebi P.

The original writing, mathematical figures, and editable writing sources in this publication are licensed under the Creative Commons Attribution 4.0 International license. You may share and adapt that material under its terms, with attribution to Komorebi P. and identification of changes. The complete license text accompanies the source package.

\noindent\url{https://creativecommons.org/licenses/by/4.0/}

The author's original Lean proofs, scripts, and other software in the accompanying release use the MIT license. Third-party software, fonts, referenced works, and other external material retain their own rights and notices. The grants above do not replace those terms. Font notices and the software license accompany the source package.

This publication preserves the accepted mathematical statements and proofs. Its source map distinguishes the formal counterparts from explanatory text, illustrations, and written proof routes. The current release report records the checks actually performed and the accepted scope of the formal library.

\vfill
\noindent\textit{Realization Calculus}
\clearpage

\pdfbookmark[0]{Contents}{visual-contents}
\begingroup
\makeatletter
\renewcommand{\@pnumwidth}{3em}
\renewcommand{\@tocrmarg}{4em}
\makeatother
\tableofcontents
\endgroup
\clearpage\phantomsection\addcontentsline{toc}{section}{How to use this guide}
\section*{How to use this guide}
Start with a picture and change one ingredient. Keep a boundary value fixed. Remove an observation. Rotate a correction direction. Tighten a solver tolerance. The following investigations show what moves, what survives, and what can fail.

Each page supplies a small declared model, a figure to inspect, and a worked change to that model. The monograph references lead to the definitions and proofs. You can follow the pictures in sequence or choose a question from the contents.

\subsection*{Three things to identify}
\begin{enumerate}
\item \textbf{The allowed objects.} Integer points, real vectors, matrices, and infinite sequences have different completion rules. A line drawn through integer markers does not create new integer solutions.
\item \textbf{The information held fixed.} An exact answer may exist but fail to preserve the prescribed data. Two pictures of exactness can describe different correction problems.
\item \textbf{The measured quantity.} A residual, a correction norm, a movement budget, and a coordinate norm are different numbers. Axis labels and metrics are part of the example.
\end{enumerate}

\subsection*{Reading the graphics}
Green marks the principal admissible object or successful route. Ochre and red distinguish comparisons, discrepancies, or excluded cases. Labels and legends identify what each color and line pattern measures. The spatial figures use genuine three-dimensional coordinates. Curves are evaluated from the formulas stated in the examples; they are not fitted to observations.

Finite windows of infinite sequences illustrate explicit formulas. An exact tail norm or a convergence condition comes from the accompanying mathematical argument, not from the end of a plotted line. Numerical drawing coordinates are rounded for display.

\subsection*{The companions}
The monograph develops the general results and their hypotheses. The Lean guide identifies exact formal declarations and their scope. These visual investigations include concrete instances and illustrative variations; they are not individually asserted to be Lean formalizations. The source package contains the plotting scripts and the recorded checks for the worked values.

\scene{01}{Read the overlaps}{Take three views of a point in $\mathbb R^3$: its $xy$, $yz$, and $xz$ coordinates. The picture follows $P=(2,1,3)$ into the three coordinate planes.}{figures/visual_companion/f01_overlapping_views.pdf}

\scenenotes{Each pairwise view omits one coordinate. The three displayed pairs agree wherever a coordinate is shared, so they describe the same point. Replacing $(2,3)$ in the $xz$ view by $(4,3)$ creates a conflict: the $xy$ view still requires $x=2$.}{Keep only the $xy$ appearance $(2,1)$. The compatible points are the whole line $(2,1,z)$. Restoring the $yz$ appearance $(1,3)$ selects $z=3$.}{Monograph Sections~1, 2}

\scene{02}{Slice the extension fiber}{Use integer variables with $k=2m$ and an independent coordinate $z$. Prescribe $k$, then look for lattice points in that slice. Only three values of $z$ are displayed.}{figures/visual_companion/f02_integer_fibers.pdf}

\scenenotes{The slice $k=2$ contains $(2,1,z)$ for every integer $z$. The slice $k=3$ misses the lattice because the required $m=3/2$ is not an integer. The free vertical coordinate changes the number of answers, while the parity condition still determines whether any answer exists.}{Prescribe $z=0$ as well as $k=2$. The circled point $(2,1,0)$ is now the unique answer. Prescribing $z=0$ cannot rescue the blocked slice $k=3$.}{Monograph Sections~3, 5}

\scene{03}{Retain the comparison that matters}{Duplicate a view and use the coefficient group $C_2$. A gauge $(a_0,a_1)$ changes the transition label by $r^{\prime}=a_0+r-a_1$. Compare $r=1$ with $r^{\prime}=0$. All arithmetic in this example is modulo two.}{figures/visual_companion/f03_boundary_comparisons.pdf}

\scenenotes{Two gauges relate the unmarked labels. The boundary marks $\phi=0$ and $\psi=1$ impose $\psi+a_0=\phi$, leaving one compatible comparison. Freezing both representative gauges instead allows only $(0,0)$, which cannot change the transition label. The three panels retain different data.}{Change the target mark to $\psi=0$ while keeping $\phi=0$. The boundary equation then selects $(0,1)$: a comparison still exists, but the chosen mark determines which one.}{Monograph Sections~15, 14}

\scene{04}{Find the missing equation}{On three ordered views, assign bit labels $c_{01}=0$, $c_{12}=0$, and $c_{02}=1$. Each two-view patch contains its own arrow, but no patch contains the full composition.}{figures/visual_companion/f04_composition_cover.pdf}

\scenenotes{The left triangle exposes the unchecked equation: the direct label is one, while the two-step label is zero. On the right, use four ordered views and all four three-view patches. Labels $c_{ij}=a_i+a_j$ from $a=(0,1,0,1)$ satisfy every face equation.}{Replace the left direct label by zero. Its composition equation then holds. This repairs that example; the pair-patch cover still provides no general guarantee that composition was checked locally.}{Monograph Sections~12, 13}

\scene{05}{Reverse the order}{Label the coordinate swap $X$ by $(1,0)$ and the second-coordinate reflection $Z$ by $(0,1)$. Start at $v=(1,2)$ and follow both orders. Matrix products act from right to left.}{figures/visual_companion/f05_order_phase.pdf}

\scenenotes{Both routes use the same two-bit labels, whose sum is $(1,1)$. Their final vectors are negatives because $ZX=-XZ$. In this model the phase $c(x,y)=x_2y_1$ records that sign. The composition law remains associative; reversing these two factors is the operation that changes the result.}{Attach a sign bit to the two-bit label. These two routes then end at different sign bits over the same label $(1,1)$, retaining the distinction that the four labels alone omit.}{Monograph Section~16}

\scene{06}{Keep the discarded detail}{Compare constant values $g_n=1$ with alternating values $f_n=(-1)^n$ on identity scalar towers. Keep only even levels. The plot displays levels zero through four and their four adjacent defects.}{figures/visual_companion/f06_hidden_detail.pdf}

\scenenotes{The sampled values agree, and each pair of alternating defects sums to zero. The coarse residual therefore cannot distinguish the candidates. The detail does: with $J(c_0,c_1)=(0,c_0,0,c_1)$, the zero coarse residual lifts to zero, leaving the entire alternating residual in $Hr=r-JCr$.}{Restore the omitted odd values. The alternating candidate has defects $-2,2,-2,2$ while the constant candidate has zero defects. Exact coherent points are unaffected by this example of cancellation in raw candidates.}{Monograph Sections~6, 10}

\scene{07}{Follow the remainder}{For the shift cascade, $Q_Nx$ sets the first $N$ coordinates to zero and leaves the tail in its original positions. Compare the constant sequence with $x_i=2^{-i}$.}{figures/visual_companion/f07_recovery_depth.pdf}

\scenenotes{The recovered region grows in both panels: every fixed coordinate eventually becomes zero in the remainder. The constant tail keeps height one, so its uniform error never decreases. The geometric tail starts at $2^{-N}$ and then decreases further, giving exactly the norm printed beneath the right panel.}{At depth $N=3$, coordinates zero, one, and two are recovered in both examples. Their infinite remainder norms are nevertheless $1$ and $1/8$, respectively.}{Monograph Section~10}

\scene{08}{Repair within the exact plane}{In Euclidean $\mathbb R^3$, take $D(x,y,z)=z$ and $C(x,y,z)=y-z$. Starting at $(2,2,3)$, first project to the locally exact plane $z=0$. The two planes encode different constraints on the same raw point.}{figures/visual_companion/f08_restricted_descent.pdf}

\scenenotes{The local correction reaches $(2,2,0)$. Solving the matching problem inside that plane reaches $(2,0,0)$, with perpendicular squared movements $9$ and $4$. Projecting instead onto the entire plane $y=z$ reaches $(2,1,1)$: matching is repaired, but the local equation has been broken.}{At the unrestricted answer, evaluate both defects: $C(2,1,1)=0$ and $D(2,1,1)=1$. The joint answer has both defects zero and total squared correction $13$.}{Monograph Section~20}

\scene{09}{How much can a preservation rule cost?}{Keep the raw point $x=(1,0,0)$ and equation $D(x)=x_1$ fixed. Vary the allowed unit correction direction $v=(\cos\theta,\sin\theta\cos\phi,\sin\theta\sin\phi)$, with protected observable $p=I-vv^{\mathsf T}$. All three coordinates use the Euclidean metric.}{figures/visual_companion/a01_protected_angles.pdf}

\scenenotes{For $0\le\theta<\pi/2$, the unique feasible correction is $a=v/\cos\theta$. Its norm is $\sec\theta$; its endpoint $x-a$ lies in the shaded exact plane. Azimuth changes the answer without changing its cost. At $\theta=\pi/2$, allowed changes cannot alter $x_1$, so this raw point becomes infeasible.}{Compare $30^\circ$ and $60^\circ$: the costs are $2/\sqrt3$ and $2$. Turning the correction toward the exact plane increases its length; every plotted line preserves its own declared observable exactly.}{Monograph Section~18}

\scene{10}{Small operators can have widely separated kernels}{Compare the row operators $D_{s,\theta}=s(\cos\theta\ \ \sin\theta)$ and $D_{s,0}=s(1\ \ 0)$, with $s>0$ and $0\le\theta\le\pi/2$. Their kernels are lines in one Euclidean plane. Sweep both scale and angle; the plot starts at $5^\circ$.}{figures/visual_companion/a02_scale_and_motion.pdf}

\scenenotes{The operator difference is $2s\sin(\theta/2)$, while the projector difference is $\sin\theta$ for the angles shown. Shrinking $s$ changes only the left panel. Each nonzero operator has least residual gain $1/s$: the scale that hides the operator difference also weakens its stability certificate.}{At $\theta=60^\circ$, reducing $s$ from $1$ to $10^{-3}$ reduces the operator difference by a thousand. The projector gap stays $\sqrt3/2$. The zero-scale operator has a different kernel and is excluded.}{Monograph Sections~18, 21}

\scene{11}{Can one residual history conceal two different outcomes?}{Let $u_n^{(p)}=\sum_{k=1}^n k^{-p}$ and $f_n^{(p)}=(s,u_n^{(p)},2^{-n})$. Protect the coordinate $s$, and measure candidate movement with the $\ell^1$ metric. Both paths start at the same point and have the same residual.}{figures/visual_companion/a03_same_residual_paths.pdf}

\scenenotes{The movement through step $N$ is exactly $\sum_{k=1}^N k^{-p}+1-2^{-N}$. For $p=2$ it converges to $1+\pi^2/6$, giving an exact limit. For $p=1$ the unprotected coordinate diverges. These conclusions use the convergent square series and divergent harmonic series, not extrapolation from the finite curves.}{After one step, both paths have residual $1/2$. Nevertheless their eventual behavior differs. The finite residual table alone cannot decide convergence: a summable bound must control movement in the candidate space.}{Monograph Section~24}

\scene{12}{Where does a family of data stop being realizable?}{Keep $Ta=(a_n/n)$ on $\ell^2$ fixed, and vary the datum $r_n^{(s)}=n^{-s}$ for $s>1/2$. The least correction satisfying its first $N$ equations has energy $E_N(s)=\sum_{n=1}^N n^{2-2s}$.}{figures/visual_companion/a04_datum_exponent_energy.pdf}

\scenenotes{An exact preimage must have coordinates $a_n=n^{1-s}$. Thus the square-series test gives realizability precisely for $s>3/2$. At $s=3/2$, the energy is harmonic and still diverges. Every finite problem remains solvable on both sides; only a uniform energy bound can pass to the Hilbert realization.}{Predict the middle curve: at $s=3/2$ its energy lies between $\log(N+1)$ and $1+\log N$. Slow growth is still unbounded growth. The plot evaluates finite sums through $N=4096$.}{Monograph Sections~26, 27, 30}

\scene{13}{When does more correction stop improving the fit?}{Let $T(a,b)=(a,b/2,0)$ and $r_h=(1,1,h)$, with Euclidean norms. The correction ball maps to a flat ellipse. Vary its radius $M$ and compare targets at heights $h=0$ and $h=1/2$.}{figures/visual_companion/a05_budget_geometry.pdf}

\scenenotes{The regularized corrections $a_\alpha=(1/(1+\alpha),2/(1+4\alpha))$ trace the exact frontier at their own budgets. At $M=\sqrt5$, the horizontal mismatch disappears; the residual floor is $|h|$. The shaded band allows data error $\delta=0.03$ and operator error $\eta=0.05$, using width $\delta+\eta M$ on one fixed correction space.}{At the marked point, $\alpha=1$, $a=(1/2,2/5)$ and $M=\sqrt{41}/10$. For $h=1/2$, the residual is $(1/2,4/5,1/2)$; its normalized direction certifies the entire budget ball.}{Monograph Sections~28, 29}

\scene{14}{Can a better measured fit give a worse answer?}{Use $T=\operatorname{diag}(1,0.02)$ and true correction $a^\dagger=(1,0)$. Add the deterministic data error $(0,0.02)$, entirely in the weak mode. The measured datum is $(1,0.02)$; all norms are Euclidean.}{figures/visual_companion/a06_regularization_modes.pdf}

\scenenotes{The exact regularized answer is $a_\alpha=(1/(1+\alpha),0.0004/(0.0004+\alpha))$. As $\alpha$ decreases, measured fit improves throughout, but the weak-mode error eventually grows. At zero-parameter limit the measured equations are fitted exactly by $(1,1)$, a distance one from the true answer. This is a finite two-mode calculation.}{At $\alpha=0.02$, the answer is $(50/51,1/51)$ and its true error is $\sqrt2/51$. Smaller measured residual alone therefore does not select the most accurate correction for this noisy datum.}{Monograph Sections~29, 32}

\scene{15}{What happens when basis vectors nearly coincide?}{Keep $T=I$, $\alpha=1$, datum $r=(0,-1)$ and proposal $a=0$ fixed in $\mathbb R^2$. Use the full basis $B_e=\left(\begin{smallmatrix}1&1\\0&e\end{smallmatrix}\right)$ for $e>0$. Its columns approach the same direction.}{figures/visual_companion/a07_basis_certificate.pdf}

\scenenotes{The actual normal residual is always $g=(0,1)$, but its coordinate covector is $q=B_e^{\mathsf T}g=(0,e)$. With $G=B_e^{\mathsf T}B_e$, the invariant square is $q^{\mathsf T}G^{-1}q=1$. At $e=0$ the columns no longer span the permitted space, so the full-space certificate formula no longer applies.}{For $e=0.01$, the raw coordinate norm is $0.01$, although the invariant norm is one. The optimizer remains $(0,-1/2)$, and the unchanged proposal still has optimizer error $1/2$.}{Monograph Sections~31, 32}

\scene{16}{Which stopping schedules preserve convergence?}{For the attainable diagonal datum $r_n=1/n^2$, choose $m=\lceil\delta^{-p/2}\rceil$, $\alpha=m^{-2}$, and solver tolerance $\varepsilon=\delta^q$, where $0<\delta\le1$. Vary the positive exponents $p,q$ instead of fixing one noise schedule.}{figures/visual_companion/a08_stopping_exponents.pdf}

\scenenotes{The proved bias, noise and solver bounds have orders $\delta^{p/4}$, $\delta^{1-p/2}$ and $\delta^{q-p}$. All vanish in the colored region $0<p<2$, $q>p$. Its color records their slowest exponent. The uncolored region loses this sufficient guarantee; it does not assert failure of every actual solver.}{The marked choice $(p,q)=(4/3,5/3)$ balances all three exponents at $1/3$. Choosing $q=p$ instead leaves a nonvanishing solver bound. The ceiling is retained in the underlying parameter rule.}{Monograph Sections~30, 32}

\end{document}
